% Slajd 3 – co prototyp już robi (kryterium: jakość i kompletność prototypu, 20 %).
\begin{frame}{Działający prototyp: aplikacja webowa na telefon i~komputer}
  \begin{columns}[T,onlytextwidth]
    \begin{column}{0.57\textwidth}
      \small
      \begin{itemize}
        \item \textbf{Katalog 6\,545 miejsc} Krakowa (OpenStreetMap) w~7~kategoriach
              i~\textbf{27 atrybutów}: schody, podjazd, winda, drzwi $\geq$~90~cm, toaleta, przewijak,
              pętla indukcyjna, miejsca odpoczynku, parking OzN w~pobliżu, cisza, pies\ldots
        \item \textbf{Profil potrzeb} zamiast pytania o~niepełnosprawność: filtruje, sortuje,
              oznacza „Dopasowane do Ciebie” i~mówi, czego brakuje.
        \item \textbf{Karta miejsca:} każda informacja ze źródłem, datą i~statusem; Aktualne / Nieaktualne /
              Popraw / Uzupełnij, pytanie do właściciela, opinie, zdjęcia, trasa dla wózka.
        \item \textbf{Mapa barier} („Yanosik dostępności”): 10 typów zgłoszeń, waga (blokuje / utrudnienie),
              głosy „nadal jest” / „naprawione”, automatyczne wygasanie.
        \item \textbf{Asystent w~języku naturalnym} (Claude z~narzędziem wyszukiwania w~bazie; bez klucza --
              reguły słów kluczowych). Pokazuje wyłącznie fakty z~bazy.
        \item \textbf{11 warstw otwartych danych}, panel właściciela, panel administratora,
              4~języki, tryb dzienny i~nocny.
      \end{itemize}
    \end{column}
    \begin{column}{0.41\textwidth}
      \screenshot{mapa-desktop}
      \vspace{0.5mm}
      \muted{Mapa z~profilem „Wózek”: miejsca (grupowane), warstwy miasta i~OSM, zgłoszone bariery
        z~wagą i~liczbą potwierdzeń. Lista obok mapy to jej tekstowa alternatywa.}
      \vspace{2mm}

      \kpi{6\,545}{miejsc z~OSM}\hfill\kpi{11}{warstw danych}\hfill\kpi{503}{testów automatycznych}
    \end{column}
  \end{columns}

  \note[item]{To nie makieta: wszystko na slajdzie działa w~repozytorium (Python stdlib + czysty JS, bez zależności).}
  \note[item]{Dane o~miejscach z~OSM: 1\,495 z~6\,545 miejsc ma już co najmniej jeden atrybut dostępności; reszta ma „brak danych” – i~tak to pokazujemy.}
  \note[item]{Asystent AI jest opcjonalny; w~kontenerze demo działa tryb reguł, więc nic nie zależy od sieci.}
\end{frame}
